
This paper documents a failed attempt to validate whether geometric properties of hidden state representations correlate with semantic entropy in large language models. The research was motivated by the computational expense of semantic entropy (~1500ms per query), which requires multi-sample generation unsuitable for production systems with sub-100ms latency requirements. We hypothesized that spectral features (participation ratio, eigenvalue decay, condition number) extracted from layers 24-31 hidden states during a single forward pass could serve as efficient uncertainty proxies. We designed an experimental pipeline to test whether these geometric features achieve Spearman correlation |ρ| > 0.4 with semantic entropy on 246 TruthfulQA test questions. The implementation comprised approximately 1,200 lines of code across five modules (data loading, hidden state extraction, geometric metrics, semantic entropy computation, correlation analysis). Data loading and model initialization succeeded, but hidden state extraction consumed over 10 hours without completion for 246 examples using Mistral-7B-v0.1, representing a 20-fold underestimate of computational requirements. No correlation measurements were obtained. The hypothesis remains untested. We quantify the computational barrier: 246 examples × 8 layers × 4096 dimensions × 2 bytes (float16) ≈ 15.4 MB tensor operations, extrapolating to approximately 12 minutes per example. This finding indicates that geometric uncertainty quantification for 7B-scale models exceeds LIGHT tier computational capacity and requires either optimized extraction infrastructure, model downsizing to sub-1B parameters, or upgraded computational resources.
